\section{方差控制与样本效率提升}
\subsection{GAE 与 temporal smoothing}
Generalized Advantage Estimation (GAE) 在 reward-to-go 和 TD error 之间平滑切换：
$$\hat{A}_t^{GAE(\gamma,\lambda)} = \sum_{l=0}^{\infty}(\gamma\lambda)^l \delta_{t+l},$$
其中 $\delta_t = r_t + \gamma V(s_{t+1}) - V(s_t)$。调节 $\lambda$ 即可在 bias 和 variance 之间做 trade-off。

\begin{importantbox}{GAE 的调节逻辑}
\begin{itemize}
    \item $\lambda=1$ 对应完全 Monte Carlo（高 variance / 低 bias）；
    \item $\lambda=0$ 对应 TD(0)（低 variance / 高 bias）；
    \item 实际上取 $\lambda \in [0.9,0.98]$ 经常能取得最优折衷。
\end{itemize}
\end{importantbox}

\subsection{Entropy 正则与 Trust Region}
为了鼓励探索，在 loss 上加入 entropy regularization：
$$\tilde{J}_{ent}(\theta) = \tilde{J}(\theta) + \beta \mathbb{E}[-\log \pi_\theta(a\mid s)].$$
Entropy term 防止策略过早确定而陷入 suboptimal mode。另一个常见做法是限制 KL divergence，形成 trust region update（如 TRPO）：
$$\max_\theta \mathbb{E}[\log \pi_\theta(a\mid s) \hat{A}] \quad \text{s.t. } D_{KL}(\pi_{\theta_{old}} \,\|\, \pi_\theta) \le \delta.$$

\begin{knowledgebox}{探索与稳定的平衡}
\begin{itemize}
    \item Entropy regularization 提供 soft constraint，鼓励概率分布更平坦；
    \item KL constraint 通过 trust region 明确限制步长，避免 policy collapse；
    \item 典型配对是 PPO（clipped objective）+ entropy bonus。
\end{itemize}
\end{knowledgebox}

\subsection{本章小结}
GAE、entropy boost 与 trust region 是在实际训练中提升稳定性和样本效率的常见手段，它们把 variance 控制在可管理范围并保持 sufficient exploration。

